\documentclass[11pt]{article}

\usepackage[a4paper,margin=1in]{geometry}
\usepackage{amsmath,amssymb,amsfonts,amsthm,mathtools,bm}
\usepackage{booktabs}
\usepackage{multirow}
\usepackage{array}
\usepackage{enumitem}
\usepackage{xcolor}
\usepackage{hyperref}
\usepackage{graphicx}
\usepackage{float}
\usepackage{setspace}
\usepackage{titlesec}

\hypersetup{
    colorlinks=true,
    linkcolor=blue,
    citecolor=blue,
    urlcolor=blue
}

\setstretch{1.08}
\setlength{\parskip}{0.45em}
\setlength{\parindent}{0pt}

\titleformat{\section}{\large\bfseries}{\thesection.}{0.5em}{}
\titleformat{\subsection}{\normalsize\bfseries}{\thesubsection.}{0.5em}{}
\titleformat{\subsubsection}{\normalsize\bfseries}{\thesubsubsection.}{0.5em}{}

\title{\textbf{LoTR: Logic-of-Thought Routing for Plug-in Reasoning}\\
\vspace{0.3em}
\large A Detailed Proposal}
\author{}
\date{}

\begin{document}

\maketitle

\section{Project Positioning}

This project studies a simple question: \emph{can reasoning be improved by dynamically matching the model's current internal state with a more suitable inference pathway at each step?}

Existing reasoning methods are often built around \emph{external program injection}: they specify a reasoning scaffold from outside the model, such as chain-of-thought prompting, self-consistency sampling, self-refinement loops, or tree/search-based exploration. These methods differ in how they organize trajectories, but they usually do not explicitly model the model's own \emph{internal state evolution} during reasoning, nor do they adapt the model's active inference pathway according to the current step-level logic demand.

Our central view is that reasoning quality depends not only on the external scaffold, but also on whether the model is using a suitable \emph{internal pathway} for the current step. At different reasoning steps, the model may be in different latent logic states. Some attention heads may support the current logic, some may be irrelevant, and some may even interfere with the current reasoning move. Therefore, a fixed pathway is suboptimal. We should instead infer the current logic state from internal information transfer, and softly route the model toward a more suitable pathway for that step.

Based on this intuition, we propose \textbf{Logic-of-Thought Routing (LoTR)}, a \emph{plug-in reasoning controller} that can be attached to diverse reasoning paradigms. LoTR does not replace CoT, SC, SR, MCTS, or other reasoning frameworks. Instead, it acts as a lightweight internal routing layer: at each sentence-level reasoning step, it reads the model's current internal state, predicts \emph{per-regime mixing weights} (sigmoid one-vs-rest in the reference code), and then softly combines precomputed logic-specific routing templates to modulate attention heads for the next step.

The method is intentionally simple. It has two main ingredients:
\begin{enumerate}[leftmargin=1.5em]
    \item a lightweight \emph{logic probe} (per layer by default, except the foundation band) that maps the current sentence-level state readout to a $K$-dimensional vector of mixing weights;
    \item a cached set of \emph{logic-specific routing templates} that determine how strongly each attention head should participate under each logic regime when that regime is active.
\end{enumerate}

The resulting method is plug-in, lightweight, reusable, and compatible with multiple reasoning paradigms.

\section{Core Thesis}

The core thesis of this project is the following:

\begin{quote}
\textbf{Reasoning should be supported by step-adaptive internal routing.} A model's current reasoning step can be associated with a latent logic state, and different logic states require different internal attention-head participation patterns. If we can identify a compact set of global logic states and their corresponding routing templates, then we can improve reasoning by dynamically selecting a more suitable pathway at each step.
\end{quote}

This thesis can be decomposed into three claims:

\begin{enumerate}[leftmargin=1.5em]
    \item \textbf{Global logic states exist.} Across diverse problems and reasoning paradigms, there exists a compact set of recurring internal logic regimes that can be discovered from reasoning trajectories.
    \item \textbf{Logic states are coupled with sparse routing patterns.} Different logic states are associated with different attention-head participation profiles, while many heads remain logic-insensitive.
    \item \textbf{Logic-conditioned routing improves plug-in reasoning.} Softly routing attention heads according to the current logic state improves reasoning quality more generally than static routing, fixed pruning, or external scaffolds alone.
\end{enumerate}

\section{Design Principles}

The method should remain clean and paradigm-like. We therefore impose the following design principles.

\subsection{Plug-in Rather than Replacement}

LoTR must be attachable to existing reasoning paradigms without redefining them. If a base method uses CoT, SC, SR, or MCTS, LoTR should simply operate on top of its stepwise generation process.

\subsection{Internal-State-Driven Rather than Rule-Driven}

We do not define logic states manually by symbolic rules, nor do we hard-code specific logic programs. Logic states are discovered from internal reasoning trajectories.

\subsection{Finite Logic Basis Rather than Instance-Level Optimization}

We do not optimize a new controller for each problem instance. Instead, we discover a compact global set of logic states offline, compute reusable routing templates for them, and perform only lightweight online combination at inference time.

\subsection{Soft Routing Rather than Hard Switching}

Different reasoning steps may exhibit mixed logic demand. Therefore, the online controller should output \emph{non-exclusive} mixing weights over logic regimes (sigmoid one-vs-rest in the implementation, allowing co-activation) and softly combine their routing templates.

\subsection{Minimal Controller Complexity}

The online controller should be simple:
\begin{itemize}[leftmargin=1.5em]
    \item no extra search,
    \item no extra planner,
    \item no separate memory module,
    \item no handcrafted fallback branch,
    \item no multi-stage controller stack.
\end{itemize}
The core online computation should ideally reduce to a small number of probe evaluations (one per non-foundation layer in the default configuration) and one weighted combination over cached templates per layer row.

\section{Problem Formulation}

Let $f_\theta$ be a frozen autoregressive language model with $L$ layers and head set
\[
\mathcal{H} = \{(\ell,h)\mid \ell=1,\dots,L,\ h=1,\dots,H_\ell\}.
\]

Consider a reasoning paradigm $\Pi$ (for example, CoT, SC, SR, or MCTS). For an input problem $x$, $\Pi$ induces a stepwise reasoning trajectory
\[
y_{1:T} = (y_1, y_2, \dots, y_T),
\]
where each $y_t$ is a sentence-level reasoning unit.

Our goal is not to redesign $\Pi$, but to augment it with a plug-in controller that modifies the model's internal routing at each step. Specifically, for each step $t$, we aim to infer \emph{logic mixing weights} over $K$ global logic regimes (implementation: one-vs-rest sigmoid outputs, \emph{not} a softmax simplex; coordinates need not sum to~$1$ and may exceed~$1$ when several regimes co-activate), and use them to construct a routing matrix
\[
W_t \in [0,1]^{L\times H},
\]
which determines the degree to which each head participates in the next-step computation (after per-entry clipping to $[0,1]$).

The resulting controller should satisfy:
\begin{enumerate}[leftmargin=1.5em]
    \item it operates at sentence level;
    \item it is independent of the outer reasoning scaffold;
    \item it reuses a fixed set of cached logic templates;
    \item it only changes attention-head participation weights.
\end{enumerate}

\section{Method Overview}

LoTR consists of two stages:

\begin{enumerate}[leftmargin=1.5em]
    \item \textbf{Offline Logic--Pathway Compilation}
    \begin{itemize}[leftmargin=1.5em]
        \item collect sentence-level reasoning states from multiple reasoning paradigms and tasks (default readout: last-layer mean-pooled hidden over each sentence span);
        \item cluster them into a compact set of global logic regimes (with soft teacher targets for probe training);
        \item train lightweight per-layer logic probes (foundation layers $0$--$2$ skipped in the default configuration);
        \item compute a routing template for each regime using the shared \texttt{HeadImportanceCalculator} path, including logic-insensitive absorption and foundation-layer forcing.
    \end{itemize}
    
    \item \textbf{Online Plug-in Routing}
    \begin{itemize}[leftmargin=1.5em]
        \item extract current sentence-span hidden states (per layer for $\ell\ge 3$ in the default stack);
        \item infer sigmoid one-vs-rest mixing weights $p_t^{(\ell)}$ (foundation rows fixed to identity gates);
        \item compose each layer row of $W_t$ by mixing $\{M_k(\ell,\cdot)\}_k$ and clipping to $[0,1]$;
        \item apply the resulting head gates on the next forward passes until the following sentence boundary.
    \end{itemize}
\end{enumerate}

The full online routing rule is still conceptually simple; the \textbf{reference implementation} instantiates it as follows.

\paragraph{Per-layer mixing weights (default).}
Let $L$ denote the number of transformer blocks and $H$ the number of attention heads per block (code uses $0$-based layer indices). The \textbf{first three layers} (indices $0,1,2$) are treated as a \emph{foundation} band: they are excluded from head-importance estimation, their cached templates are set to all ones, and their online gates are fixed to~$1$ (analogous to absorbing low-accuracy early-layer probes into an always-on backbone band). For each remaining layer $\ell\in\{3,\dots,L-1\}$, we read a sentence-span mean-pooled hidden vector $r_t^{(\ell)}\in\mathbb{R}^d$ from that layer's residual stream (stacked hidden states from the backbone), apply a \emph{layer-specific} linear probe, and set
\[
p_t^{(\ell)} = \sigma\!\left(W^{(\ell)} r_t^{(\ell)} + b^{(\ell)}\right)\in (0,1)^K,
\]
where $\sigma$ is applied coordinate-wise (one-vs-rest). The layer row of the gate matrix is
\[
W_t(\ell,\cdot) = \mathrm{clip}\!\left(\sum_{k=1}^{K} p_t^{(\ell)}(k)\, M_k(\ell,\cdot),\, 0,\, 1\right)\in[0,1]^H,
\]
and $W_t(0\!:\!2,\cdot)$ is identically~$1$.

\paragraph{Legacy single-probe path (optional checkpoint).}
Some artifacts use one shared probe on the \emph{last-layer} sentence vector $r_t^{(L-1)}$ only; then a single vector $p_t=\sigma(W r_t+b)$ is used for every layer row $W_t(\ell,\cdot)$ with the same mixing formula (foundation rows may still be forced to~$1$ if recorded in metadata).

\paragraph{Objects.}
\begin{itemize}[leftmargin=1.5em]
    \item $r_t^{(\ell)}$ (or $r_t$ in the legacy path) is the current sentence-level state readout;
    \item $p_t^{(\ell)}(k)$ (or $p_t(k)$) is the \emph{sigmoid} activation for regime~$k$;
    \item $M_k\in[0,1]^{L\times H}$ is the cached routing template for regime~$k$.
\end{itemize}

\section{Offline Stage: Global Logic-State Discovery}

\subsection{Step 1: Build a Joint Reasoning-State Pool}

We collect sentence-level reasoning trajectories from diverse reasoning paradigms and tasks. The purpose is not to compare paradigms, but to build a broad and heterogeneous pool of reasoning states from which global logic states can be discovered.

Suppose we have a set of reasoning paradigms
\[
\mathcal{P} = \{\Pi_1,\Pi_2,\dots,\Pi_m\}
\]
and a set of tasks
\[
\mathcal{D} = \{\mathcal{D}_1,\mathcal{D}_2,\dots,\mathcal{D}_n\}.
\]
For each paradigm--task pair, we run the frozen model and extract sentence-level states from generated reasoning trajectories. This yields a pooled collection
\[
\mathcal{R} = \{r_t^{(i)}\}_{i,t},
\]
where each $r_t^{(i)}$ is a sentence-level internal state representation.

This joint pool is important because our hypothesis is that logic is \emph{problem-free}: global logic states are not tied to a specific task or outer reasoning scaffold. Different paradigms may express different subsets of the global logic space, but the space itself should be shared.

\subsection{Step 2: Sentence-Level State Readout}

For each sentence $y_t$, we extract a compact internal state vector
\[
r_t = \Phi(y_t; f_\theta)\in\mathbb{R}^d.
\]

The reference pooling pipeline for \emph{clustering and probe training} uses the \textbf{last} hidden block's token representations (mean over the current sentence span), i.e.
\[
r_t = \frac{1}{|y_t|}\sum_{j\in y_t} h_j^{(L-1)} \quad\text{(0-based indexing)}.
\]
At inference, the default additionally forms per-layer readouts $\{r_t^{(\ell)}\}_{\ell\ge 3}$ from earlier blocks using the stacked hidden-state tuple, while the foundation band remains gated to identity.

Optionally, this state can be augmented by a small number of trajectory-dynamic statistics, but the main method should remain valid even with a simple sentence-level readout.

The point of $r_t$ is not to summarize task semantics alone, but to capture the model's current internal reasoning condition at that step.

\subsection{Step 3: Global Logic Clustering}

We cluster the pooled reasoning states $\mathcal{R}$ into $K$ clusters:
\[
\mathcal{C} = \{C_1,C_2,\dots,C_K\}.
\]

Each cluster corresponds to a global logic state:
\[
z_k \leftrightarrow C_k.
\]

The logic states are intended to be:
\begin{itemize}[leftmargin=1.5em]
    \item \textbf{compact}: small enough to cache and interpret;
    \item \textbf{recurrent}: shared across tasks and paradigms;
    \item \textbf{distinct}: associated with different routing patterns;
    \item \textbf{softly composable}: training uses soft targets; at test time sigmoid outputs allow multi-regime co-activation without a simplex constraint.
\end{itemize}

We can determine $K$ using a cluster quality criterion such as silhouette score, or by balancing separability and coverage in the pooled reasoning-state space. In the paper, this should be described as a compact global logic basis rather than over-emphasizing any single clustering heuristic.

\paragraph{Reference implementation (clustering loop).}
The repository runs $K$-means for $K$ in a bounded range and selects the configuration with the best silhouette score on the pooled features $\{r_t\}$. Pseudo-labels $\hat{z}_t$ are the hard cluster assignments.

\paragraph{Soft teacher targets for probe training.}
In addition to hard labels, each $r_t$ receives a \emph{soft} target vector $y_t\in(0,1)^K$ by applying a Gaussian kernel in squared Euclidean distance to centroids,
\[
y^{\text{soft}}_{t,k} \propto \exp\!\left(-\frac{\|r_t-c_k\|_2^2}{2\tau^2}\right),
\]
normalizing over $k$, then blending with the one-hot hard label: $y_t = 0.85\, y^{\text{soft}} + 0.15\, e_{\hat{z}_t}$ (renormalized). This stabilizes probe training when clusters overlap.

\subsection{Step 4: Logic Labels as Regimes, Not Rules}

Each cluster defines a \emph{logic regime}. We may assign interpretive names in the analysis section, such as decomposition-like, verification-like, elimination-like, correction-like, or synthesis-like. However, these names should be treated as descriptive summaries rather than hard supervision labels.

Thus, the method does not claim that logic is explicitly given by symbolic rules. Instead, it claims that recurring internal reasoning regimes can be discovered and then used as control units.

\subsection{Repository Tooling (Data Integrity Defaults)}

The reference repository provides scripts to (i) collect sentence-level trajectories with frozen backbones, (ii) compile artifacts (clustering $\rightarrow$ soft targets $\rightarrow$ probe training $\rightarrow$ template estimation), and (iii) run plug-in evaluation. \textbf{Correctness filtering:} trajectory rows may be labeled with task-aware grading (same extractor and scorer as the MoT evaluation stack); the compiler supports filtering to \texttt{sample\_correct=true} rows so incorrect trajectories do not distort probe or template estimation. \textbf{Diagnostics:} online runs emit structured diagnostics under \texttt{generation\_meta} by default (sigmoid summaries, gate histograms, rollups) to audit whether regime readouts separate and whether gates stay in a healthy dynamic range (many heads near identity, a minority modulated).

\section{Offline Stage: Logic Probe}

\subsection{Pseudo Labels from Clustering}

After clustering, each pooled state $r_t$ receives a pseudo label
\[
\hat{z}_t \in \{1,\dots,K\}.
\]

These pseudo labels, together with the soft targets $y_t$ from the previous subsection, supervise the logic probe.

\subsection{Probe Architecture (Default: Per-Layer, Foundation Band)}

The reference implementation maintains $L$ small linear probes $\{(W^{(\ell)},b^{(\ell)})\}_{\ell=0}^{L-1}$ (one per transformer block). \textbf{Layers $\ell\in\{0,1,2\}$ are not trained}: empirical probe accuracy is unreliable in very early layers, so those layers are excluded from probe training and treated as a \emph{foundation band} (see also routing templates and online gates below). For $\ell\ge 3$, each probe is trained against the \emph{same} pooled sentence vectors $\{r_t\}$ as the clustering features (last-layer mean-pooled hidden states in the default data pipeline), using the \emph{same} soft targets $y_t$---a deliberate engineering choice until per-layer readouts are stored in trajectories.

\subsection{Probe Outputs and Training Objective}

For each trained layer $\ell\ge 3$, logits are $z^{(\ell)}_t = W^{(\ell)} r_t + b^{(\ell)}\in\mathbb{R}^K$, and online mixing weights use the \textbf{sigmoid} (one-vs-rest) map
\[
p_t^{(\ell)} = \sigma(z^{(\ell)}_t)\in(0,1)^K.
\]

This matches the \texttt{slm\_agent} probe convention (BCE / sigmoid rather than softmax) and allows multiple regimes to co-activate.

The probe is trained by minimizing \textbf{binary cross-entropy with logits} against the soft vector targets $y_t$:
\[
\mathcal{L}_{\text{probe}}
=
\frac{1}{N}\sum_{t=1}^{N}\mathrm{BCEWithLogits}\!\left(z^{(\ell)}_t,\, y_t\right).
\]

This is \emph{not} softmax cross-entropy on hard labels alone; the soft teacher carries distance-to-centroid information.

\section{Offline Stage: Logic-Specific Routing Templates}

\subsection{Motivation}

The probe alone only tells us which logic state is likely active. We also need to know which internal pathway best supports that logic state. This is the role of the routing template cache.

For each logic state $k$, we want to compute a routing template
\[
M_k \in [0,1]^{L\times H},
\]
where $M_k(\ell,h)$ indicates how strongly head $(\ell,h)$ should participate when the model is in logic state $k$.

\subsection{Head Relevance Under a Logic Regime (Importance Cache)}

Let cluster $k$ index a set of \emph{input} examples used only for offline estimation (token sequences sampled from trajectories assigned to cluster~$k$). For each layer $\ell$ and head $h$, the repository ships a \textbf{vendored} \texttt{HeadImportanceCalculator} inside the LoTR codebase (no sibling-project import at runtime): it extracts a per-head contribution tensor along the residual stream, averages over the sequence dimension to obtain a vector $w_{\ell,h}\in\mathbb{R}^d$, applies the \emph{layer-$\ell$} linear probe, and maps logits through a \textbf{sigmoid} to obtain $u=\sigma(W^{(\ell)} w_{\ell,h}+b^{(\ell)})\in(0,1)^K$. With one-hot axes aligned to regimes, the scalar score recorded for regime~$k$ is the coordinate $u_k$. Averaging over samples in cluster~$k$ yields a raw importance matrix $I_{\ell,h}^{(k)}$.

\paragraph{Important implementation detail.}
The \textbf{foundation band} $\ell\in\{0,1,2\}$ skips this computation entirely in code: importance rows are left at zero and are later overwritten by the template normalization rules plus an explicit ``all ones'' assignment for those layers (see below).

\paragraph{Relation to the conceptual energy formula.}
The conceptual ``axis-aligned energy'' story in earlier drafts is \emph{superseded in code} by the vendored calculator above (probe projection + sigmoid + coordinate selection + averaging), aligned with the historical \texttt{slm\_agent} reference implementation.

\subsection{Template Construction}

After computing the raw tensor $I_{\ell,h}^{(k)}$, the repository applies two post-processing steps:
\begin{enumerate}[leftmargin=1.5em]
    \item \textbf{Logic-insensitive absorption.} For each $(\ell,h)$, if the empirical standard deviation of $\{I_{\ell,h}^{(k)}\}_{k=1}^{K}$ falls below a small threshold, the head is declared logic-insensitive and all $M_k(\ell,h)$ are set to~$1$ (no separate whitelist file; the head is absorbed into every regime template as fully on).
    \item \textbf{Per-layer min--max normalization.} For each fixed $k$, each layer row $M_k(\ell,\cdot)$ is min--max normalized into $[0,1]$ across heads; degenerate rows become all ones.
\end{enumerate}

Finally, for the foundation band, the implementation \textbf{forces} $M_k(\ell,h)=1$ for all $k$ and all heads whenever $\ell\in\{0,1,2\}$, guaranteeing identity gates on those rows regardless of probe noise.

The resulting templates $M_k\in[0,1]^{L\times H}$ are cached in the artifact and reused at inference time.

\subsection{Logic-Insensitive Heads (Cross-Reference)}

The operational definition of logic-insensitive heads ($\sigma_{\ell,h}$ small $\Rightarrow$ all $M_k(\ell,h)=1$) is implemented exactly as described in \textbf{Template Construction} above; there is no separate whitelist JSON. The foundation band is an additional, layer-wise ``always on'' mechanism independent of~$\sigma_{\ell,h}$.

\section{Online Stage: Plug-in Logic Routing}

At inference time, LoTR is attached to a base reasoning paradigm $\Pi$.

\subsection{Step 1: Read Current Sentence-Level State(s)}

After the model completes a sentence-ending token (delimiter-based segmentation in the reference code), we extract sentence-span mean-pooled hidden states. \textbf{Default:} for each layer $\ell\ge 3$, $r_t^{(\ell)}$ is computed from that layer's hidden state tensor; the backbone returns the full hidden-state tuple from HuggingFace (stored by reference without per-token extra stacking) and stacks it \emph{once per sentence boundary} for efficiency.

\subsection{Step 2: Predict Regime Mixing Weights}

For each layer $\ell\ge 3$, the corresponding probe outputs logits $z_t^{(\ell)}$ and mixing weights $p_t^{(\ell)}=\sigma(z_t^{(\ell)})$ coordinate-wise. Foundation rows do not run probes: $W_t(0\!:\!2,\cdot)$ remains at~$1$.

\subsection{Step 3: Compose Routing Template}

For each layer row separately,
\[
W_t(\ell,h)=\mathrm{clip}\!\left(\sum_{k=1}^{K} p_t^{(\ell)}(k)\, M_k(\ell,h),\,0,\,1\right).
\]
This is a \textbf{linear mixture} of regime templates followed by clipping---not a Boolean conjunction. When many heads are logic-irrelevant, their template rows are near one and the composed gates remain near one (identity-like behavior), while a smaller set of logic-sensitive heads can be driven toward lower gates when the active mixture down-weights their supporting regimes.

\subsection{Step 4: Modulate Attention Heads}

For each head $(\ell,h)$, let $o_{\ell,h}^{(t)}$ denote its attention output slice at the current forward (implementation: head outputs are reshaped, multiplied by the scalar gate $W_t(\ell,h)$, and reshaped back). No other part of the model is changed. The outer reasoning paradigm remains exactly the same.

\subsection{Interpretation}

This mechanism can be read as follows:
\begin{itemize}[leftmargin=1.5em]
    \item the probes identify which logic regimes are currently active \emph{per layer row} (default) or globally (legacy checkpoint);
    \item the cached templates encode offline-estimated head--regime alignment;
    \item the weighted composition \emph{down-weights} heads whose templates are small under the currently active mixture; it is an associative cache, not an online gradient-based proof of ``negative utility.''
\end{itemize}

\subsection{Online Diagnostics (Default-On in Code)}

Each sentence-boundary update records rich diagnostics inside \texttt{generation\_meta} (including per-layer sigmoid summaries, gate sparsity/fractiles, and rollups). Diagnostics are \textbf{enabled by default} (\texttt{LOTR\_STEP\_DIAGNOSTICS=1}); set \texttt{LOTR\_STEP\_DIAGNOSTICS=0} to disable for benchmarking. Trajectory scripts persist the full \texttt{generation} dict (and \texttt{lotr\_steps}) for offline analysis.

\section{Why the Method Is Clean}

A central strength of LoTR is that it remains structurally simple.

\subsection{Only Two Core Objects}

The entire method is built around two reusable objects:
\begin{enumerate}[leftmargin=1.5em]
    \item a logic probe,
    \item a cache of logic-specific routing templates.
\end{enumerate}

\subsection{Online Computation Is Minimal}

The online routing reduces to (default) one small linear probe forward per non-foundation layer at each sentence boundary, plus the per-row mixing and clipping that yields $W_t$. No search, no planner, no extra branch scoring, and no new objective are introduced at inference time. The dominant runtime cost remains the frozen LM forward; the backbone already requests full hidden-state tuples when needed, and the reference implementation stacks them only at sentence boundaries to avoid per-token redundant copies.

\subsection{No Redesign of the Base Reasoning Framework}

LoTR does not redefine CoT, SC, SR, or MCTS. It only routes the internal pathway more appropriately at each step.

\subsection{No Handcrafted Logic Rules}

The logic states are discovered from trajectories rather than specified by a symbolic rule set.

\subsection{Finite and Cacheable Control Units}

The clustering step is not for interpretability alone. It is what allows continuous reasoning space to be compressed into a finite cacheable control basis.

\section{Expected Advantages}

We expect LoTR to provide the following benefits.

\subsection{Better Step-Level Reasoning Alignment}

Different reasoning steps require different internal pathway configurations. LoTR adapts pathway participation to the current step rather than relying on a fixed global pattern.

\subsection{Compatibility Across Reasoning Paradigms}

Since it is plug-in, LoTR can be attached to multiple external reasoning frameworks.

\subsection{Suppression of Logic Interference}

If some heads are irrelevant or disruptive under the current logic state, LoTR can down-weight them through the composed routing matrix.

\subsection{Natural Sparsity}

Because many heads are logic-insensitive and many logic distributions are expected to be sharp, the effective routing pattern should be relatively sparse.

\subsection{Interpretability}

The discovered logic states and their corresponding routing templates provide a natural analysis interface:
\begin{itemize}[leftmargin=1.5em]
    \item which logic regimes recur across paradigms?
    \item which heads are logic-sensitive?
    \item which heads are shared backbones?
    \item which routing patterns help which reasoning steps?
\end{itemize}

\section{Main Experimental Plan}

\subsection{Part I: Plug-in Effectiveness Across Reasoning Paradigms}

\textbf{Goal.} Show that LoTR works as a general plug-in controller.

\textbf{Base paradigms.}
\begin{itemize}[leftmargin=1.5em]
    \item Vanilla
    \item CoT
    \item Self-Consistency (SC)
    \item Self-Refine (SR)
    \item Plan-and-Solve
    \item MCTS-style reasoning
\end{itemize}

\textbf{Comparison.}
For each paradigm, compare:
\begin{itemize}[leftmargin=1.5em]
    \item base paradigm,
    \item base paradigm + LoTR.
\end{itemize}

\textbf{Metrics.}
\begin{itemize}[leftmargin=1.5em]
    \item accuracy / exact match,
    \item pass@1 or recall where appropriate,
    \item average token usage,
    \item latency,
    \item cost-adjusted performance,
    \item robustness under more difficult settings.
\end{itemize}

\textbf{Target conclusion.}
LoTR should improve multiple reasoning paradigms without requiring paradigm-specific redesign.

\subsection{Part II: Logic-State Discovery Analysis}

\textbf{Goal.} Show that the discovered clusters are meaningful as recurring reasoning regimes.

\textbf{Analyses.}
\begin{itemize}[leftmargin=1.5em]
    \item visualization of the logic-state space,
    \item distribution of different paradigms over logic states,
    \item distribution of tasks over logic states,
    \item qualitative interpretation of each cluster.
\end{itemize}

\textbf{Target conclusion.}
The discovered logic states are neither arbitrary nor paradigm-specific; they recur across diverse trajectories.

\subsection{Part III: Routing Specificity}

\textbf{Goal.} Show that logic-specific templates matter.

\textbf{Key experiments.}
\begin{itemize}[leftmargin=1.5em]
    \item matched template vs random template,
    \item matched template vs average template,
    \item off-diagonal swap matrix: use template $M_j$ on steps classified as logic $i$,
    \item shuffled cluster labels.
\end{itemize}

\textbf{Target conclusion.}
The gain does not come from generic reweighting alone, but from matching the current logic state to the right routing template.

\subsection{Part IV: Not Just Static Pruning or Surface Features}

\textbf{Goal.} Distinguish LoTR from simpler alternatives.

\textbf{Baselines.}
\begin{itemize}[leftmargin=1.5em]
    \item static head pruning,
    \item average head importance,
    \item activation-only routing,
    \item routing by sentence length,
    \item routing by step index,
    \item routing by task ID,
    \item routing by paradigm ID.
\end{itemize}

\textbf{Target conclusion.}
LoTR should outperform baselines that only exploit static salience or surface metadata.

\subsection{Part V: Transferability}

\textbf{Goal.} Support the claim that logic is problem-free.

\textbf{Experiments.}
\begin{itemize}[leftmargin=1.5em]
    \item train logic states on a subset of paradigms, test on unseen paradigms,
    \item train on a subset of tasks, test on unseen tasks,
    \item optionally test across nearby backbones.
\end{itemize}

\textbf{Target conclusion.}
LoTR captures reusable logic regimes rather than overfitting to one task or one scaffold.

\subsection{Part VI: Sparsity and Efficiency}

\textbf{Goal.} Validate the intended routing structure.

\textbf{Metrics.}
\begin{itemize}[leftmargin=1.5em]
    \item distribution of sigmoid outputs (including coordinate maxima, sums, and an optional entropy of the normalized diagnostic view),
    \item proportion of logic-sensitive vs logic-insensitive heads,
    \item overlap between routing templates,
    \item number of strongly suppressed heads per step,
    \item online overhead.
\end{itemize}

\textbf{Target conclusion.}
The method is lightweight, sparse in practice, and structurally consistent with its design hypothesis.

\section{Critical Ablations}

The following ablations are particularly important.

\subsection{No Clustering}

Replace the discrete logic basis with a fully continuous controller.

\textbf{Purpose.} Show that a compact finite logic basis is useful for cacheability, interpretability, and stable routing.

\subsection{Softmax vs.\ Sigmoid Probe (Routing Specificity)}

The reference implementation uses \textbf{sigmoid one-vs-rest} probes (BCE training). An ablation may reintroduce a softmax probe + cross-entropy training to test whether forced simplex competition changes routing specificity or end-task accuracy.

\textbf{Purpose.} Quantify the empirical difference between mutually-exclusive versus co-activatable regime readouts under the same template cache.

\subsection{No Logic-Insensitive Absorption}

Force all heads to be logic-dependent.

\textbf{Purpose.} Show that preserving invariant backbone heads is beneficial.

\subsection{Average Template Only}

Use a single global template for all steps.

\textbf{Purpose.} Show that dynamic logic-conditioned routing is the key, not just global reweighting.

\subsection{Random or Surface-Level Clustering}

Cluster using random assignments or obvious surface features.

\textbf{Purpose.} Show that the quality of the logic partition is central.

\section{Potential Risks and How to Address Them}

\subsection{Risk 1: Clusters May Reflect Surface Style Rather Than Logic}

This is a legitimate concern. The intended response is not to claim perfect semantic logic labeling, but to show that the discovered regimes are functionally meaningful through:
\begin{itemize}[leftmargin=1.5em]
    \item recurrence across paradigms,
    \item recurrence across tasks,
    \item matched-routing superiority,
    \item superiority over surface-feature baselines.
\end{itemize}

\subsection{Risk 2: Gains May Be Interpreted as Generic Dynamic Pruning}

This should be addressed by demonstrating:
\begin{itemize}[leftmargin=1.5em]
    \item LoTR operates at sentence-level logic adaptation,
    \item matched routing is better than random or average routing,
    \item static pruning baselines do not recover the same gains,
    \item routing is driven by internal state rather than fixed scenario metadata.
\end{itemize}

\subsection{Risk 3: Pseudo Labels from Clustering May Be Noisy}

This is true in principle, but the method does not depend on exact semantic labels. It depends on whether the partition yields reusable and functionally distinct routing templates. This is testable through the routing-specificity experiments.

\subsection{Risk 4: The Controller May Add Complexity Without Clear Benefit}

This is why the paper should keep the online controller maximally simple and push all complexity into offline compilation. The clean \emph{per-layer-row} formula (with clipping),
\[
W_t(\ell,\cdot)=\mathrm{clip}\!\left(\sum_k p_t^{(\ell)}(k)\,M_k(\ell,\cdot),\,0,\,1\right),
\]
remains the conceptual center of the method (specialized to a single global $p_t$ only in legacy checkpoints).

\section{Main Contribution Summary}

If successful, this paper would contribute the following:

\begin{enumerate}[leftmargin=1.5em]
    \item \textbf{A new plug-in reasoning controller.}  
    LoTR introduces step-level logic-conditioned routing without replacing existing reasoning paradigms.
    
    \item \textbf{A compact global logic basis.}  
    It proposes to discover a finite set of recurring logic states from heterogeneous reasoning trajectories.
    
    \item \textbf{A simple logic--pathway coupling mechanism.}  
    Each logic state is associated with a reusable routing template over attention heads.
    
    \item \textbf{A new view of reasoning enhancement.}  
    The work shifts emphasis from external reasoning scaffolds alone to matching the model's current internal condition with a more suitable inference pathway.
\end{enumerate}

\section{Concise One-Paragraph Abstract Draft}

We propose \textbf{Logic-of-Thought Routing (LoTR)}, a plug-in controller for inference-time reasoning that dynamically adjusts the model's internal pathway according to its current sentence-level logic regime mixture. Existing reasoning methods often rely on externally injected programs, such as prompting strategies, refinement loops, or search procedures, but they do not explicitly adapt the model's internal computation to the current reasoning step. LoTR addresses this gap by discovering a compact set of global logic regimes from heterogeneous reasoning trajectories, training lightweight per-layer probes (with an early-layer foundation band) to predict \emph{sigmoid} one-vs-rest mixing weights from internal readouts, and softly combining cached logic-specific routing templates---computed with an in-repository head-importance calculator (vendored; no external project path required at runtime)---to modulate attention heads. The method is simple, reusable, and compatible with diverse reasoning paradigms including CoT, self-consistency, self-refinement, and search-based reasoning. We hypothesize that different regimes are coupled with different sparse attention-head participation patterns, and that routing heads according to the current mixture can down-weight heads that are misaligned with the active regimes while leaving logic-insensitive and foundation heads near identity. If validated, LoTR would introduce a clean new interface for plug-in reasoning: not changing the outer reasoning scaffold, but improving reasoning by matching each step to a more suitable internal pathway.

\section{Conclusion}

LoTR is designed to be a clean and modular reasoning proposal. It does not attempt to hand-code logic or replace existing reasoning methods. Instead, it introduces a minimal internal routing layer that reads the model's current reasoning state(s), maps them to a compact global logic space, and activates a more suitable pathway for the next sentence span. The main promise of this work is not complexity, but simplicity: a bank of lightweight probes (with a small early-layer foundation band), a finite cache of routing templates produced by shared head-importance machinery, and a stepwise soft routing rule that can be attached to many reasoning paradigms---together with default diagnostic telemetry for scientific scrutiny. If the empirical evidence supports the central hypotheses, LoTR could provide a new and elegant direction for plug-in reasoning.

\end{document}